\thispagestyle{plain}
\begin{center}
    \Large \textbf{\uppercase{Abstract}}
\end{center}

\vspace{3\baselineskip}

\noindent
The development of high-performance solid-state batteries (SSBs) is a critical step towards next-generation energy storage, but the discovery of novel solid electrolytes with high ionic conductivity remains a significant bottleneck. Traditional experimental and first-principle computational methods are often resource-intensive and time-consuming. This thesis explores the application of machine learning (ML) to accelerate the screening and discovery of promising solid electrolyte candidates. Utilizing the comprehensive Database of Doped Solid Electrolytes (DDSE) \cite{DDSE2021}, we develop and evaluate a predictive model for ionic conductivity. The model is built upon a feature set derived exclusively from fundamental elemental properties, intentionally avoiding the need for complex crystal structure information. This structure-agnostic approach enables rapid preliminary screening of a vast chemical space. Various ML algorithms were assessed, with a Gradient Boosting Regressor demonstrating the highest predictive accuracy, achieving a coefficient of determination (\(R^2\)) of 0.75 and a Root Mean Square Error (RMSE) of 0.51 log(S/cm) on the held-out test set. The results validate the efficacy of using elemental-property-based features to forecast ionic conductivity, establishing a computationally efficient framework to identify and prioritize materials for further experimental synthesis and validation in the pursuit of advanced solid-state electrolytes.


\vspace{\baselineskip}

\noindent
\textbf{Keywords}: Solid-State Batteries, Ionic Conductivity, Machine Learning, Materials Informatics, Feature Engineering, Gradient Boosting
